
\subsubsection{具体分析}

问题一要求判断单根金属A圆柱体在给定坐标下能否使微构体导通，属于三维几何建模与距离判定问题。本问需建立考虑周期性边界条件的圆柱体-电极距离计算模型。

首先建立圆柱体的几何表示，将金属A表示为以两端点连线为轴、半径为30nm的直圆柱体，接着根据周期性边界条件将超出微构体边界的部分按规则平移回立方体内部，然后计算圆柱体表面到两侧电极面（X=-5000和X=+5000）的最短距离，若两侧距离均不超过隧穿距离1.8nm则判定导通。具体分析流程如图\ref{fig:analysis1_flow}所示。

\begin{figure}[ht]
  \begin{center}
    \begin{tikzpicture}[x=1cm, y=0.7cm]
      \node[box] (n11) at (0, 0) {读取坐标};
      \node[box] (n12) at (5, 0) {周期裁剪};
      \node[box] (n13) at (10, 0) {分段生成};
      \node[box] (n22) at (5, -3) {距离计算};
      \node[box] (n21) at (0, -3) {左电极距};
      \node[box] (n23) at (10, -3) {右电极距};
      \node[box] (n31) at (0, -6) {导通判断};
      \node[box] (n32) at (5, -6) {概率统计};
      \node[box] (n33) at (10, -6) {结果输出};
      \draw[arrow] (n11) -- (n12);
      \draw[arrow] (n12) -- (n13);
      \draw[arrow] (n23) -- (n22);
      \draw[arrow] (n22) -- (n21);
      \draw[arrow] (n13) -- ++(0,-1.0) -| (n23);
      \draw[arrow] (n21) -- (n31);
      \draw[arrow] (n31) -- (n32);
      \draw[arrow] (n32) -- (n33);
    \end{tikzpicture}
    \caption{问题一分析流程图}
    \label{fig:analysis1_flow}
  \end{center}
\end{figure}

\subsubsection{模型准备}

在建立数学模型之前，需要明确周期性边界条件的处理方式和金属A圆柱体的几何参数。

金属A为直圆柱体，其轴线由两端点$\mathbf{p}_1=(x_1,y_1,z_1)$和$\mathbf{p}_2=(x_2,y_2,z_2)$确定，底面半径$R_A=30$nm，长度$L_A=5000$nm。附件1包含12个金属A颗粒的端点坐标数据，每个颗粒构成一个独立的微构体配置。数据的X坐标范围为$[-5000, 3660.69]$至$[-3939.16, 5000]$，Y坐标范围为$[-500, 273.51]$至$[-472.32, 420.23]$，Z坐标范围为$[-500, 441.67]$至$[-413.21, 500]$，所有坐标均在微构体范围$[-5000, 5000]^3$内或其边界上。

周期性边界条件的处理规则为：当导电介质的任一部分超出微构体边界时，将超出部分沿越界方向的反方向平移一个微构体边长（10000nm），使其由相对侧边界重新进入微构体。若平移后仍在其他方向超出边界，则重复上述操作直至所有部分均回到微构体范围内。该规则等价于三维环面拓扑结构。

该问题本质上是一个三维空间中圆柱体到平面的最短距离计算问题，由于周期性边界条件和仅有一个颗粒，关键在于确定圆柱体是否能同时接触两侧电极面。综上所述，本节完成了数据理解和问题定位，接下来将建立几何模型并进行求解。
